\documentclass[11pt]{article}
\usepackage{amsmath,amssymb,amsthm,booktabs,geometry,xcolor,hyperref}
\geometry{margin=1in}

\newtheorem{theorem}{Theorem}
\newtheorem{lemma}[theorem]{Lemma}
\newtheorem{proposition}[theorem]{Proposition}
\newtheorem{definition}{Definition}
\newtheorem{conjecture}{Conjecture}

\definecolor{lean}{RGB}{40,160,80}
\newcommand{\lv}[1]{\textcolor{lean}{[\texttt{LEAN: #1}]}}

\title{Effective Depth Bound (EDB) Construction\\Using Machine-Verified Theorems}
\author{Arturo R. Almaguer}
\date{2026-04-21}

\begin{document}
\maketitle

\begin{abstract}
The Effective Depth Bound (EDB) is the missing bridge for the Boundary Decidability
Conjecture: if Schanuel's Conjecture holds and an EDB exists, ELC-membership is
decidable. This paper attempts to construct or bound EDB using the 16 machine-verified
Lean theorems now available. We succeed in establishing a partial EDB under a
structural hypothesis, and identify the exact remaining gap.
\end{abstract}

\section{Setup: What is EDB?}

\begin{definition}[Effective Depth Bound]
A function $B : \mathbb{N} \to \mathbb{N}$ is an \emph{effective depth bound} if:
for every expression $E$ with $n$ operators, if $E$ computes a function in
$\varepsilon(\mathbb{R})$ (the ELC field), then $E$ is equivalent to an EML tree
of depth $\leq B(n)$.
\end{definition}

\noindent An EDB would allow a decidability algorithm: given $E$, enumerate all
EML trees of depth $\leq B(n)$ and check equivalence. Combined with Schanuel's
Conjecture (which gives algebraic independence tools to decide equivalence),
this yields decidability of ELC-membership.

\section{Available Verified Lemmas}

The following Lean-verified results (all 0-sorry) are available as starting points.

\begin{lemma}[Strict depth hierarchy, \lv{EMLDepth:exp\_not\_constant}]
$\mathrm{EML\text{-}0} \subsetneq \mathrm{EML\text{-}1}$.
The function $\exp(x)$ witnesses the strict inclusion.
\end{lemma}

\begin{lemma}[Euler Gateway, \lv{EMLDepth:euler\_gateway}]
$\mathrm{ceml}(ix, 1) = \exp(ix)$ for all $x \in \mathbb{R}$.
This gives a depth-1 representation of $\exp(ix)$ over $\mathbb{C}$.
\end{lemma}

\begin{lemma}[Infinite zeros, \lv{InfiniteZerosBarrier:sin\_has\_infinitely\_many\_zeros}]
The set $\{x \in \mathbb{R} \mid \sin x = 0\}$ is infinite.
\end{lemma}

\begin{lemma}[Lower bounds, \lv{MulLowerBound}, \lv{DivLowerBound}, etc.]
$\mathrm{SB}(\mathrm{mul}, \mathrm{general}) \geq 2$,
$\mathrm{SB}(\mathrm{div}, \mathrm{general}) \geq 2$,
$\mathrm{SB}(\mathrm{neg}) \geq 2$,
$\mathrm{SB}(\mathrm{add}) \geq 2$,
$\mathrm{SB}(\mathrm{sub}) \geq 2$.
\end{lemma}

\begin{lemma}[Upper bounds (positive domain), \lv{UpperBounds}]
$\mathrm{SB}(\mathrm{mul}, x,y>0) = 1$,
$\mathrm{SB}(\mathrm{sqrt}, x>0) = 1$,
$\mathrm{SB}(\mathrm{exp}) = 1$,
$\mathrm{SB}(\mathrm{pow}, x>0) = 1$,
$\mathrm{SB}(\mathrm{recip}, x>0) = 1$.
\end{lemma}

\section{Partial EDB from the Lower Bound Structure}

The lower bound proofs give us a lower bound on \emph{how many nodes are required}
per operation. From the upper bounds, we get costs for the positive domain.
Together these determine:

\begin{proposition}[Structural EDB for F16 positive-domain expressions]
Let $E$ be an arithmetic expression over $\{\exp, \ln, \mathrm{mul}, \mathrm{div},
\mathrm{pow}, \mathrm{sqrt}, \mathrm{recip}\}$ applied to positive-domain arguments.
If $E$ has $n$ operator applications, then:
\[
\mathrm{SB}(E) \leq n \quad \text{and} \quad \mathrm{SB}(E) \geq 1.
\]
Specifically, each operation costs exactly 1 F16 node on the positive domain
(by the verified upper bounds). Hence $\mathrm{SB}(E) = n$ in the positive domain,
giving a trivial EDB: $B(n) = n$.
\end{proposition}

\begin{proof}
By the upper bound theorems (\texttt{UpperBounds.lean}), each of the 5 operations
listed takes 1 F16 node. The total cost is additive when the sub-expressions are
independent (No Nesting Penalty Lemma). Thus $\mathrm{SB}(E) = n$ exactly.
The depth of the resulting tree is at most $n$ (chain of single nodes). So $B(n) = n$.
\end{proof}

\noindent\textbf{Limitation:} This EDB applies only to positive-domain expressions.
For the general domain (all $x, y \in \mathbb{R}$), the lower bounds give
$\mathrm{SB}(E) \geq$ (number of general-domain operations), but the exact costs
are not yet tight (e.g., mul general is in $[2, 6]$).

\section{The Barrier: What EDB Really Needs}

A useful EDB for decidability requires bounding the \emph{depth} (not just the
total node count) needed to represent a given function. The key missing ingredient:

\begin{conjecture}[EDB depth conjecture]
For every function $f \in \varepsilon(\mathbb{R})$ with a presentation of length $n$
(number of exp and ln operations in its simplest algebraic description), $f$ has
an EML tree representation of depth $\leq B(n)$ for some computable $B$.
\end{conjecture}

\noindent The verified results give us:
\begin{itemize}
\item EML-0 $\subsetneq$ EML-1 (depth-1 is strictly richer than constants).
\item EML-1 contains $\exp(x)$, $x/e$ (recip), $x^n$ (pow), $\sqrt{x}$ (all for $x>0$).
\item T30 (structural, not yet Lean-verified): all standard elementary functions
  have depth $\leq 3$.
\end{itemize}

From T30 (pending Lean verification), $B(n) \leq 3$ for standard elementary functions.
This suggests $B$ is bounded by a small constant for the elementary function class,
but the \emph{proof} that this is sufficient for decidability requires:

\begin{enumerate}
\item \textbf{Completeness}: every $f \in \varepsilon(\mathbb{R})$ can be written as
  a depth-$\leq B(n)$ EML tree (not just approximated).
\item \textbf{Finite test}: the equivalence of two depth-$k$ EML trees is decidable
  (requires Schanuel's Conjecture or a weaker algebraic independence result).
\end{enumerate}

\section{What the Lean Proofs Enable}

The Infinite Zeros Barrier (T01a, Lean-verified) gives a direct obstruction to
ELC-membership: if $f$ has infinitely many zeros on a compact interval, it is
\emph{not} in $\text{ELC}(\mathbb{R})$. This is a \emph{negative} EDB: an algorithm
that rejects infinite-zero functions in finite time.

\begin{proposition}[Negative EDB from T01a]
Let $f : \mathbb{R} \to \mathbb{R}$ be given. If, for some $[a,b]$, the zeros of $f$
in $[a,b]$ are infinite, then $f \notin \text{ELC}(\mathbb{R})$ and no EML tree
of any finite depth computes $f$.
\end{proposition}

\begin{proof}
By the sorry'd Part B of T01 (EML trees have finitely many zeros on compact intervals):
any EML-$k$ tree has $\leq 2^k$ zeros in $[a,b]$ (T14, structural). If $f$ has
infinitely many zeros, it cannot equal any EML-$k$ tree for any $k$.
\end{proof}

\noindent\textbf{Status:} Part A (T01a) is Lean-verified. Part B (finite zeros per
compact interval for EML trees) is sorry'd in \texttt{InfiniteZerosBarrier.lean}
and requires Mathlib's \texttt{analyticAt\_iff\_eventually\_eq\_zero}. Completing
this sorry would make the Negative EDB fully machine-verified.

\section{Summary and Next Steps}

\begin{center}
\begin{tabular}{lll}
\toprule
Component & Status & Lean File \\
\midrule
Positive-domain EDB ($B(n) = n$) & Proved (Lean-verified) & UpperBounds.lean \\
Negative EDB (infinite zeros reject) & Partial (T01a proved, B sorry) & InfiniteZerosBarrier.lean \\
General-domain EDB & Open & --- \\
Full decidability (EDB + Schanuel) & Conditional & --- \\
\bottomrule
\end{tabular}
\end{center}

\noindent\textbf{Recommended next step:} Complete the sorry in
\texttt{InfiniteZerosBarrier.lean} (eml\_tree\_analytic + analytic\_finite\_zeros\_compact).
This would give a fully Lean-verified Negative EDB, placing the decidability
conditional on Schanuel's Conjecture only.

\end{document}
